\documentclass[11pt,a4paper]{article}

\usepackage{url}
\usepackage[utf8]{inputenc}
\usepackage[english]{babel}
\usepackage{amsmath}
\usepackage{graphicx}
\usepackage{geometry}
\usepackage{booktabs}
\usepackage{tabularx}
\usepackage[table,dvipsnames]{xcolor}
\geometry{margin=2cm}
\usepackage{soul}
\usepackage{subcaption}
\usepackage{hyperref}

% Let more floats share a page so large tables are not deferred to the end.
\renewcommand{\topfraction}{0.9}
\renewcommand{\bottomfraction}{0.8}
\renewcommand{\textfraction}{0.07}
\renewcommand{\floatpagefraction}{0.75}
\setcounter{topnumber}{3}
\setcounter{totalnumber}{4}

\title{Project Report Group 27: Comparing GA, PSO, and ACO for\\
Pathfinding in 2D Mazes}
\author{\large Manon Leijser \and Nemo Ingendaa \and Robert Blaauwendraad}
\date{\today}

\hypersetup{
    colorlinks=true,
    linkcolor=cyan,
    citecolor=cyan,
    urlcolor=black
}

\begin{document}

    \maketitle

    \section{Introduction}
    Classical pathfinding algorithms such as A* are fast and optimal in static, fully known grids, but they degrade in two situations that arise constantly in practice. When the environment hides a deceptive dead end on the apparent route to the goal, a greedy distance heuristic is led astray; and when the environment changes after a route has been computed, a single-agent solver must discard its work and search again. Population-based natural computing algorithms offer a different proposition: rather than committing to one solution, a population maintains many candidate routes at once and can, in principle, adapt collectively when the landscape shifts beneath it.

    Comparative studies of these algorithms exist but stop short of the questions that matter for adaptation. Shrestha et al.~\cite{shrestha2024} benchmark a Genetic Algorithm (GA), Particle Swarm Optimization (PSO), and Ant Colony Optimization (ACO) on a shared 2D environment, and Sariff and Buniyamin~\cite{sariff2009} compare GA and ACO across static maps of increasing complexity. Both judge the algorithms by their endpoint results, namely whether the goal is reached and how long the final path is, and both hold the environment fixed for the duration of a run. Neither reports how the population distributes itself across the search space while solving, and neither tests what happens when an established route is invalidated mid-run. This leaves a concrete gap: the existing literature treats GA, PSO, and ACO as interchangeable metaheuristics ranked by a final score, when the more informative distinction is mechanistic. The three algorithms differ fundamentally in how their agents share information, and that difference should govern both how a population explores and how it recovers from disruption. Endpoint metrics on static maps cannot reveal this, which is precisely why it is worth measuring.

    The information-sharing mechanism is the axis this project isolates. A GA is \emph{decoupled}: active agents exchange nothing during evaluation, and knowledge passes only vertically between generations through selection and recombination. PSO is \emph{tightly coupled}: every particle immediately sees the swarm's global best and is pulled toward it. ACO is \emph{environmentally coupled}: agents communicate indirectly through pheromone deposited on the grid, so the environment itself acts as a shared, decaying memory. These three regimes span a clean spectrum of shared memory, from none, through continuous and global, to distributed and environmental. Pathfinding is an ideal testbed for studying that spectrum because the search space is controlled and fully observable, a deceptive or non-stationary structure can be engineered deliberately, and ACO's spatial pheromone maps directly onto the navigation space without an interpolation layer that would confound the comparison. Conclusions about how shared memory helps or hurts adaptation then transfer to any population-based optimiser facing deceptive or changing objectives.

    This motivates the central research question:

    \begin{quote}
        \emph{How does the degree of agent coupling that distinguishes GA, PSO, and ACO shape the way a population adapts when its environment becomes deceptive, non-stationary, or multimodal?}
    \end{quote}

    \noindent Three sub-questions make this concrete. SQ1 concerns architecture: how does each algorithm's coupling mechanism determine the way its population distributes itself across the search space over time? SQ2 concerns the environment: how does the type of environmental challenge interact with each architecture to produce distinct failure and recovery patterns? SQ3 concerns scale: to what extent does population size modulate adaptive capacity, and does increasing scale compensate for or amplify architectural weaknesses?

    \paragraph{Expected behaviour.}
    Three predictions, tied to the sub-questions, were fixed before the confirmatory run so that the results could confirm or refute them rather than be read selectively afterwards. The first (H1, architecture) is that each regime produces a distinct entropy signature: PSO should show the sharpest collapse on meeting a strong attractor, since the global best draws all particles to one position; GA should remain the most stable, since decoupled agents cannot lock onto a single misleading path together; and ACO should fall between them, declining gradually as pheromone builds and erodes. The second (H2, environment) is that deception stresses PSO most, because tight coupling accelerates commitment to the misleading attractor; non-stationarity stresses ACO most, because pheromone on invalidated paths persists until evaporation allows rerouting; and multimodality also stresses PSO most in spatial spread, because social pressure forces the swarm onto one route even when alternatives are equally good. The third (H3, scale) is that larger populations amplify PSO's weakness rather than compensate for it, benefit GA most because decoupled variation improves with more independent samples, and leave ACO comparatively insensitive because pheromone dynamics dominate over ant count.

    \section{Related Work}
    Robot path planning has long served as a benchmark for population-based search. Sariff and Buniyamin~\cite{sariff2009} applied GA and ACO to static environments of varying complexity and found ACO consistently faster than GA, reaching the goal in fewer iterations across every tested complexity and with parameters easier to adapt as complexity grew. Shrestha et al.~\cite{shrestha2024} extended this comparison to all three algorithms on a common 2D environment with parameters tuned for a waypoint representation, in which a path is a sequence of random coordinate points scattered across an open space. That sparse representation demands different navigation dynamics from a dense grid. Both studies compare endpoint performance on a fixed map; neither characterises how the population behaves internally while solving, nor what happens when the environment changes mid-run.

    Measuring internal behaviour requires a diversity signal. Masisi et al.~\cite{masisi2008} used Shannon entropy to quantify the structural diversity of a classifier ensemble and found that greater diversity accompanied higher ensemble accuracy. Their setting is classifier ensembles rather than path planning, but the entropy-based diversity measure transfers directly: applied to the spatial distribution of agents across the maze, it gives a continuous indicator of premature convergence and recovery speed rather than a single final score. Ansari et al.~\cite{ansari2015} studied efficiency in complex pathfinding and observed that existing search algorithms, A* among them, can require a large number of iterations to find a path; that efficiency framing motivates the iteration budgets used here, while the present work targets the underlying causes of those delays.

    Against this background, the contribution is fourfold. The first extension introduces a \textbf{dynamic environment} by blocking an established route mid-run (Maze~2, Section~\ref{environments}), testing adaptation under disruption, which Shrestha et al.\ explicitly leave as future work. The second measures \textbf{population-internal dynamics} through position-wise Shannon entropy rather than endpoint metrics alone, giving a continuous record of spatial diversity and premature convergence. The third \textbf{isolates individual environmental properties}, namely deception, non-stationarity, and multimodality, using purpose-built maze structures instead of varying difficulty monotonically through random obstacle counts. The fourth \textbf{frames the algorithms as distinct information architectures} and derives failure-mode predictions from that framing rather than treating them as interchangeable. Together these extensions address how memory architectures degrade and recover when the environment itself changes.

    \section{Approach}
    The simulation framework is built from scratch in Python, reusing the core GA and PSO routines from earlier course assignments and adapting them to grid-based navigation. All three algorithms are population-based optimisers searching the same space of start-to-goal paths under the same budget, and they differ only in how agents share information. GA agents never communicate within a generation, and information survives only through selected offspring; PSO agents all see one global best and are drawn toward it each step; ACO agents deposit and read pheromone on the grid, a shared memory that decays over time. Holding the problem, the budget, and the population fixed lets behavioural differences be attributed to this coupling axis rather than to incidental implementation choices.

    \subsection{Environment and grid model}
    The environment is a 2D grid in which agents navigate from a start node to a goal node in as few steps as possible. A discrete grid was chosen deliberately over a continuous space so that ACO's spatial memory maps directly onto the navigation space, isolating the architectural differences without an interpolation layer that would itself need tuning. Walls block movement, and four-neighbour adjacency means any valid path is a connected sequence of cells, a property the GA crossover relies on.

    \subsection{Algorithms and parameter choices}
    Because the objective is to compare information architectures rather than to maximise absolute performance, parameters were not tuned per maze. Per-maze tuning would over-fit and mask the very architectural differences under study; instead, parameters are fixed from literature sources chosen to match each algorithm's most comparable problem formulation (Table~\ref{tab:parameters}). Since no single paper tunes all three on a grid-based maze with cell-by-cell navigation, sources are combined and each adaptation is documented.

    \begin{table}[ht]
        \centering
        \small
        \begin{tabularx}{\textwidth}{lllX}
\toprule
\textbf{Algorithm} & \textbf{Parameter} & \textbf{Value} & \textbf{Source / rationale} \\
\midrule
\textbf{GA}
& Encoding & Cell sequence (var.\ length) & Path is the natural representation; preserves validity \cite{benhlima2018,tu2003} \\
& Selection & Roulette wheel & Follows Shrestha et al. \cite{shrestha2024} \\
& Crossover & Two-point at common cells (0.5) & Shared-cell splice keeps paths connected \\
& Mutation & Truncate-and-regrow (0.3) & Rate from \cite{shrestha2024}; operator adapted to keep paths valid \\
\midrule
\textbf{PSO}
& Inertia $\omega$ & 1.0 & Directly from Shrestha et al. \cite{shrestha2024} \\
& Cognitive $c_1$ & 0.1 & As above \\
& Social $c_2$ & 0.2 & As above ($r_1,r_2\sim U(0,1)$ are runtime draws, not parameters) \\
\midrule
\textbf{ACO}
& Pheromone exp.\ $\alpha$ & 1 & Grid-based ACO of Akka and Khaber \cite{akka2018} \\
& Heuristic exp.\ $\beta$ & 5 & As above \\
& Deposit $Q$ & 2 & As above \\
& Evaporation $\rho$ & 0.1 & From Shrestha et al. \cite{shrestha2024}; long pheromone memory \\
& Initial pheromone & 0.8 & As above \\
\bottomrule
        \end{tabularx}
        \caption{Fixed algorithm parameters and the source or rationale for each. Population size is not listed here because it is a deliberate experimental factor (Section~\ref{design}), tested at 20, 50, and 150 for all three algorithms.}
        \label{tab:parameters}
    \end{table}

    The GA encodes each chromosome as a variable-length sequence of cells representing the path from the start node. An initial binary direction encoding, a fixed-length string of up/down/left/right moves amenable to standard bit-flip mutation, produced a 0\% success rate across all trials. Pilot sweeps confirmed the failure was structural rather than a matter of tuning or compute: it persisted across a range of mutation rates and chromosome lengths and at budgets up to 50{,}000 generations. The cause is that direction-string crossover preserves no inheritable invariant, because splicing two valid move sequences almost never yields a valid, continuous path. The cell-sequence representation of Lamini et al.~\cite{benhlima2018} and the variable-length grid chromosome of Tu and Yang~\cite{tu2003} resolve this, and differ from the fixed-width waypoint chromosome of Shrestha et al.~\cite{shrestha2024}, which encodes points scattered in open Euclidean space. In a confined grid, the continuous path is itself the natural representation. Roulette-wheel selection and the 0.3 mutation rate are retained from Shrestha et al., but the genetic operators are adapted to the encoding: two-point crossover is restricted to cells shared between parents so that the child path stays connected, and mutation cuts the path at a random point and regrows a new valid route to the goal rather than flipping individual cells.

    PSO adopts its parameters directly from Shrestha et al.~\cite{shrestha2024}. Particle positions are projected onto the grid, and the velocity update blends inertia, the personal best, and the global best. ACO follows the grid-based, cell-by-cell formulation of Akka and Khaber~\cite{akka2018}, which matches the present setup more closely than the feasible-node graph of Sariff and Buniyamin~\cite{sariff2009} or the waypoints of Shrestha et al.; its evaporation and initial pheromone follow Shrestha et al. Colony size is abstracted away from these baselines because population scale is itself a tested factor. Rather than the disparate per-algorithm sizes reported by Shrestha et al.\ (GA 70--100, PSO 30--40, ACO 20--30 across their scenarios), population is aligned across all three algorithms at 20, 50, and 150. Comparing architectures at identical scales separates coupling effects from a mere population advantage and treats scale as a deliberate axis (SQ3) rather than a confound.

    \subsection{Maze environments} \label{environments}
    Three maze types are used, each engineered to isolate one environmental challenge so that a failure can be attributed to a specific property rather than to generic difficulty. For each type, 10 distinct instances are generated that preserve the structural property of that type, with 10 independent trials per instance. This yields 100 runs per condition and separates algorithmic stochasticity (trial seed) from maze-instance variance (instance seed).

    The first maze, the Shortest Path Trap, isolates deception. A dead end sits on the straight line to the goal, so agents drawn toward the goal enter it, and escaping requires moving away from the goal against the distance heuristic. By construction the maze is perfect: exactly one route reaches the goal, with the trap corridor and a separate dead wedge as the only dead-end branches, so any agent that reaches the goal does so by the unique and therefore optimal path. This is the cleanest test of whether an architecture can override a misleading gradient.

    The second maze, the Sudden Wall, isolates non-stationarity. Its base is perfect, with one extra opening added as a shortcut so that two routes exist initially, a short one and a longer alternative. The shortcut is walled at iteration $T=100$, after the population has settled into it, leaving the long route as the unique path. Because the only thing that differs from a static maze is this disruption, any change in behaviour after $T=100$ is attributable to how each architecture handles invalidated memory.

    The third maze, Parallel Paths, isolates multimodality. Two routes of different length lead from the start, a shorter one and a longer one via a small detour, with no third route. It is the only maze with more than one solution path, and hence the only one on which path optimality can exceed $1$. It tests whether a population can spread across multiple optima or whether coupling forces premature commitment to one.

    \subsection{Experimental design} \label{design}
    A single consolidated three-factor experiment crosses algorithm (GA, PSO, ACO; SQ1), maze type (Shortest Path Trap, Sudden Wall, Parallel Paths; SQ2), and population size (20, 50, 150; SQ3), rather than running three separate studies. The interesting questions, namely whether the effect of scale depends on the environment and whether coupling scales differently, are interaction effects that a set of one-factor-at-a-time studies cannot recover. The population range spans the 20--100 reported by Shrestha et al.\ and extends to 150. Each of the 27 cells is evaluated on 10 maze instances with 10 trials each, giving 100 runs per cell and 2{,}700 runs in total. Random seeds are controlled so that all three algorithms face the exact same generated instances, and the iteration limit is set to $10\times$ the A* shortest-path length for each instance, so the budget scales with intrinsic difficulty and is identical across algorithms on a given instance.

    \subsection{Evaluation metrics} \label{metrics}
    Because the research question concerns adaptive behaviour rather than mere goal-reaching, internal-dynamics metrics are paired with grounding performance metrics, and each is linked to the sub-question it serves.

    Spatial diversity addresses SQ1, namely how each architecture's coupling determines the way its population distributes itself across the search space over time. It operationalises the exploration--exploitation balance as how widely agents spread across structurally distinct regions at a given iteration, and is quantified by position-wise Shannon entropy
    \[
        H_i = -\sum_{j \in \Sigma} f_{ij}\,\log f_{ij},
    \]
    where $f_{ij}$ is the fraction of agents in cell $j$ at iteration $i$. Entropy is the natural quantity here because it distinguishes the three coupling regimes directly: a tightly coupled swarm concentrates mass on few cells (low $H$), while an exploratory population spreads it out (high $H$). The shape of the curve is the main diagnostic, with a sharp cliff indicating rapid collapse, a slow drift indicating lagging environmental memory, and a noisy plateau indicating stochastic wandering. Three scalars summarise it for statistical comparison: the \emph{time to 50\% entropy loss}, measured from the peak so that collapse is captured only after maximum spread; the \emph{diversity floor}, the minimum entropy after the peak; and the \emph{mean entropy} over the run.

    Adaptation time addresses SQ2, namely how an environmental challenge, here non-stationarity, interacts with each architecture to produce distinct recovery patterns. It counts the iterations from the disruption at $T=100$ until mean entropy first returns to 80\% of its pre-disruption value. A threshold is needed because entropy recovers gradually and ``recovered'' has no natural cutoff; the same 80\% threshold is applied to every algorithm, size, and instance for fairness. Runs that never recross 80\% before the limit are excluded from the adaptation average and captured instead by the success rate, so a slow recoverer is not silently scored as fast.

    Two grounding metrics complete the picture. Success rate, the fraction of the 100 runs per cell that reach the goal within the iteration limit, anchors the entropy analysis in actual task completion, since high diversity is worthless if the goal is never found. Path optimality, the final path length divided by the A* optimum for successful runs only, is informative only on Parallel Paths, the one maze with more than one solution route; there it separates algorithms that settle for the slightly longer detour from those that find the shorter route. On the Shortest Path Trap and the Sudden Wall the solution is unique, so every successful run is optimal by construction and optimality adds nothing beyond the success rate.

    The statistical procedure was fixed before the confirmatory run. Each cell reports a mean and a 95\% bootstrap confidence interval (10{,}000 resamples) for success rate, iterations, optimality, and the entropy scalars, and entropy curves carry shaded 95\% bands across the 100 runs. Algorithm contrasts within a maze$\times$size cell use the Mann--Whitney $U$ test with Holm correction, since success and optimality are bounded and non-normal, and are reported with rank-biserial effect sizes so that differences are judged by magnitude as well as significance given the large sample. SQ3, on whether and how population size modulates adaptive capacity, is tested by fitting each metric with algorithm, maze, and population as factors and examining the population main effect and the algorithm$\times$population interaction (Table~\ref{tab:factoreffects}). Full per-cell intervals and the complete table of pairwise tests accompany the code.

    \section{Results}
    All numbers below come from the full confirmatory run of 2{,}700 trials (100 per cell), summarised in Table~\ref{tab:summary}. The three sub-questions are addressed in turn.

    \begin{table}[ht]
        \centering
        \small
        \rowcolors{2}{gray!12}{white}
        \begin{tabular}{llrrrrrrrr}
            \toprule
            Algo & Maze & Pop & SR & Iters & Opt & $t_{50\%}$ & Floor & $\bar{H}$ & Adapt \\
            \midrule
            GA & SPT & 20 & \textbf{0.00} & 2304 & N/A & 10 & 0.00 & 0.39 & N/A \\
            GA & SPT & 50 & \textbf{0.00} & 2304 & N/A & 10 & 0.07 & 0.53 & N/A \\
            GA & SPT & 150 & \textbf{0.00} & 2304 & N/A & 10 & 0.28 & 0.67 & N/A \\
            GA & SW & 20 & \textbf{0.42} & 3070 & 1.00 & 67 & 0.82 & 2.19 & 8 \\
            GA & SW & 50 & \textbf{0.47} & 2681 & 1.00 & 361 & 1.95 & 2.90 & 7 \\
            GA & SW & 150 & \textbf{0.65} & 2125 & 1.00 & 32 & 2.41 & 3.37 & 7 \\
            GA & PP & 20 & 1.00 & 1 & 1.01 & N/A & 3.33 & 3.33 & N/A \\
            GA & PP & 50 & 1.00 & 1 & 1.00 & N/A & 3.87 & 3.87 & N/A \\
            GA & PP & 150 & 1.00 & 1 & 1.00 & N/A & 4.24 & 4.24 & N/A \\
            \midrule
            PSO & SPT & 20 & 1.00 & 491 & 1.00 & 417 & 3.75 & 3.31 & N/A \\
            PSO & SPT & 50 & 1.00 & 406 & 1.00 & N/A & 5.11 & 3.85 & N/A \\
            PSO & SPT & 150 & 1.00 & 339 & 1.00 & N/A & 6.00 & 4.25 & N/A \\
            PSO & SW & 20 & 1.00 & 670 & 1.00 & 280 & 3.90 & 3.18 & 5 \\
            PSO & SW & 50 & 1.00 & 624 & 1.00 & 220 & 5.07 & 3.76 & 4 \\
            PSO & SW & 150 & 1.00 & 586 & 1.00 & 230 & 6.04 & 4.28 & 6 \\
            PSO & PP & 20 & 0.97 & 108 & 1.02 & 14 & \textbf{1.41} & 1.06 & N/A \\
            PSO & PP & 50 & 1.00 & 85 & 1.02 & 18 & \textbf{1.56} & 1.10 & N/A \\
            PSO & PP & 150 & 1.00 & 85 & 1.01 & 18 & \textbf{1.82} & 1.23 & N/A \\
            \midrule
            ACO & SPT & 20 & 1.00 & 351 & 1.00 & N/A & 4.05 & 3.56 & N/A \\
            ACO & SPT & 50 & 1.00 & 298 & 1.00 & N/A & 5.17 & 4.24 & N/A \\
            ACO & SPT & 150 & 1.00 & 276 & 1.00 & N/A & 6.29 & 4.83 & N/A \\
            ACO & SW & 20 & 1.00 & 532 & 1.00 & 190 & 3.61 & 3.74 & \textbf{1} \\
            ACO & SW & 50 & 1.00 & 491 & 1.00 & 160 & 4.76 & 4.58 & \textbf{1} \\
            ACO & SW & 150 & 1.00 & 464 & 1.00 & N/A & 5.72 & 5.28 & \textbf{1} \\
            ACO & PP & 20 & 1.00 & 79 & 1.01 & N/A & 3.38 & 2.15 & N/A \\
            ACO & PP & 50 & 1.00 & 78 & 1.00 & N/A & 4.00 & 2.43 & N/A \\
            ACO & PP & 150 & 1.00 & 78 & 1.00 & N/A & 4.48 & 2.62 & N/A \\
            \bottomrule
        \end{tabular}
        \caption{Summary metrics per experimental cell (3 algorithms $\times$ 3 maze types $\times$ 3 population sizes; 100 runs each). Bold marks the headline results: GA's complete failure and partial recovery (SR column), PSO's diversity collapse on Parallel Paths (lowest Floor), and ACO's near-instant Sudden-Wall recovery (Adapt). \textbf{Algo}=algorithm; \textbf{Maze} (SPT=Shortest Path Trap, SW=Sudden Wall, PP=Parallel Paths); \textbf{Pop}=population size; \textbf{SR}=success rate ($0.00$--$1.00$); \textbf{Iters}=mean iterations to goal (or to the limit if unsolved); \textbf{Opt}=mean path optimality ($\geq 1.00$, $1.00$=A* optimum); \textbf{$t_{50\%}$}=iterations from the entropy peak to 50\% loss; \textbf{Floor}=minimum post-peak entropy (bits); \textbf{$\bar{H}$}=mean entropy (bits); \textbf{Adapt}=iterations to recover 80\% entropy after the disruption. \texttt{N/A} means the metric does not apply: Opt when SR$=0$, Adapt outside Sudden Wall, and $t_{50\%}$ when entropy never falls to half its peak.}
        \label{tab:summary}
    \end{table}

    \subsection{Diversity dynamics}
    The entropy curves with 95\% bands in Figure~\ref{fig:entropy} answer SQ1, on how coupling shapes the spatial distribution of the population. On the deceptive trap (a) and the Sudden Wall (b) the two coupled architectures behave alike: ACO and PSO both climb to a high diversity plateau of roughly $5$ bits, with ACO reaching it faster so that its mean entropy is marginally the higher of the two (Table~\ref{tab:summary}). GA behaves very differently across these two mazes. On the trap it collapses to a floor near $0.5$ bits and never recovers, whereas on the wall it holds a steady, moderate level around $3$ bits.

    Parallel Paths (c) inverts the picture. PSO collapses onto a single route and sits at the lowest diversity, the lowest floor in Table~\ref{tab:summary}, while ACO spreads across both routes and rises highest. GA cannot be compared on equal terms here because it reaches the goal in its first generation, so its tabulated entropy reflects only the initial random spread of one generation rather than a sustained trajectory.

    The clearest reading of H1 is therefore the opposite of its prediction. Wherever GA actually searches over time, on the trap and the wall, it is the least diverse architecture rather than the most. These contrasts are statistically unambiguous: on the trap and the wall GA's diversity floor falls below both coupled architectures with the maximal rank-biserial effect ($r \approx -1.0$, Holm-corrected $p < 10^{-30}$), and on Parallel Paths PSO's floor sits below ACO's ($r \approx -0.8$, $p < 10^{-21}$), the signature of the multimodal collapse.

    \begin{figure}[ht]
        \centering
        \begin{subfigure}{0.32\textwidth}
            \includegraphics[width=\textwidth]{imgs_project/entropy_shortest_path_trap_pop50.png}
            \caption{Shortest Path Trap}
        \end{subfigure}\hfill
        \begin{subfigure}{0.32\textwidth}
            \includegraphics[width=\textwidth]{imgs_project/entropy_sudden_wall_pop50.png}
            \caption{Sudden Wall}
        \end{subfigure}\hfill
        \begin{subfigure}{0.32\textwidth}
            \includegraphics[width=\textwidth]{imgs_project/entropy_parallel_paths_pop50.png}
            \caption{Parallel Paths}
        \end{subfigure}
        \caption{Population diversity (Shannon entropy, bits) over iterations at population 50, with 95\% bands across 100 runs. On the deceptive trap (a) and the Sudden Wall (b), ACO and PSO both climb to a high plateau near $5$ bits, ACO faster; GA collapses to about $0.5$ bits on the trap but holds a moderate $\sim 3$ bits on the wall. The dashed line in (b) marks the disruption at $T=100$. On Parallel Paths (c), PSO sits lowest having collapsed onto a single route, while ACO rises highest; GA is not visible because it solves the maze in its first generation, leaving a single high-entropy point at the left edge rather than a curve.}
        \label{fig:entropy}
    \end{figure}

    \subsection{Task performance across environments}
    Success rate by population (Figure~\ref{fig:success}) addresses SQ2, on how each environmental challenge interacts with each architecture. GA solves Parallel Paths completely, fails outright on the deceptive trap at every size, and only partially solves the Sudden Wall, where its success climbs from $0.42$ to $0.65$ as population grows. The coupled architectures solve all three mazes, ACO at 100\% everywhere and PSO at 100\% on the trap and wall and $0.97$--$1.00$ on Parallel Paths. The gap is significant wherever GA struggles: against both coupled architectures it shows the maximal effect on the trap ($r = -1.0$, Holm-corrected $p < 10^{-43}$ at every size) and a smaller but significant gap on the wall that narrows as GA catches up ($r = -0.58$ at population 20 falling to $-0.35$ at 150).

    Optimality discriminates only on Parallel Paths, where the table shows PSO most often settling for the slightly longer detour (Opt up to $1.02$) while GA and ACO almost always find the shorter route (Opt $\approx 1.00$). GA's success on that maze should be read with care, since it reaches the goal in a single generation (Iters $=1$): the maze has no deceptive structure, so even random initial paths already connect start to goal. A raw success rate would otherwise credit GA with an ability it never demonstrates on the harder mazes.

    Adaptation after the Sudden-Wall disruption (Figure~\ref{fig:adapt}) contradicts H2. ACO, predicted to be slowest because of stale pheromone, instead recovers its diversity in about one iteration at every size, and PSO recovers within four to six, because both maintain a high diversity floor and have therefore already explored the alternative route before the wall appears. GA, the architecture expected to adapt fastest, is the slowest to recover, at seven to eight iterations, and converts that recovery into a successful solve only $42$--$65\%$ of the time.

    \begin{figure}[ht]
        \centering
        \begin{subfigure}{0.32\textwidth}
            \includegraphics[width=\textwidth]{imgs_project/success_rate_shortest_path_trap.png}
            \caption{Shortest Path Trap}
        \end{subfigure}\hfill
        \begin{subfigure}{0.32\textwidth}
            \includegraphics[width=\textwidth]{imgs_project/success_rate_sudden_wall.png}
            \caption{Sudden Wall}
        \end{subfigure}\hfill
        \begin{subfigure}{0.32\textwidth}
            \includegraphics[width=\textwidth]{imgs_project/success_rate_parallel_paths.png}
            \caption{Parallel Paths}
        \end{subfigure}
        \caption{Success rate by population size and maze. GA solves Parallel Paths, fails the deceptive trap entirely, and partially solves the Sudden Wall with success rising as population grows; the coupled architectures PSO and ACO solve all three mazes.}
        \label{fig:success}
    \end{figure}

    \begin{figure}[ht]
        \centering
        \includegraphics[width=0.55\textwidth]{imgs_project/adaptation_time_sudden_wall.png}
        \caption{Adaptation time on the Sudden Wall (iterations to recover 80\% of pre-disruption entropy). ACO recovers almost immediately and PSO within a few iterations; GA is the slowest to recover its diversity and still reaches the goal only $42$--$65\%$ of the time.}
        \label{fig:adapt}
    \end{figure}

    \subsection{Effect of population scale}
    Scale, the subject of SQ3, raises the diversity floor monotonically across all three algorithms and mazes as population grows from 20 to 50 to 150. Larger populations therefore help every architecture maintain spatial spread rather than amplifying a weakness in PSO as H3 predicted, and for PSO and ACO the higher floor translates into fewer iterations to the goal (Table~\ref{tab:summary}).

    GA benefits most visibly where it has room to. On the Sudden Wall its success climbs from $0.42$ to $0.65$ as the larger population raises its floor and lets more independent lineages stumble onto the open route. On the deceptive trap, however, scale is not enough: the floor rises but never crosses the threshold needed to escape, and success stays at $0.00$ for every size.

    A balanced factorial analysis makes this precise (Table~\ref{tab:factoreffects}). The population main effect is significant for every metric, and the algorithm$\times$population interaction is significant for all but optimality. This is the formal statement of H3's outcome: scale helps, but by an architecture-dependent amount, rather than uniformly amplifying PSO's weakness.

\begin{table}[ht]
    \centering
    \small
    \begin{tabular}{lrrrr}
        \toprule
         & \multicolumn{2}{c}{Population main effect} & \multicolumn{2}{c}{Algorithm $\times$ population} \\
        \cmidrule(lr){2-3}\cmidrule(lr){4-5}
        Metric & $F$ & $p$ & $F$ & $p$ \\
        \midrule
        Success rate      & $7.0$    & $9\times10^{-4}$  & $5.4$  & $3\times10^{-4}$ \\
        Iterations        & $9.4$    & $9\times10^{-5}$  & $3.0$  & $0.018$ \\
        Diversity floor   & $1008.0$ & $<10^{-250}$      & $34.7$ & $3\times10^{-28}$ \\
        Mean entropy      & $723.0$  & $<10^{-250}$      & $12.1$ & $1\times10^{-9}$ \\
        Optimality (PP)   & $20.1$   & $3\times10^{-9}$  & $1.2$  & $0.33$ \\
        \bottomrule
    \end{tabular}
    \caption{Factor effects from a balanced factorial analysis (Type II ANOVA) over all 2{,}700 runs. The population main effect is significant for every metric; the algorithm$\times$population interaction is significant for all but optimality, confirming that larger populations help by an architecture-dependent amount (SQ3). Success rate is binary and iterations are censored at the budget, so for those two the $F$ statistic is read alongside the rank tests rather than on its own. Optimality is restricted to Parallel Paths, the only maze on which it varies.}
    \label{tab:factoreffects}
\end{table}

    \section{Discussion}
    The headline finding is largely a reversal of the hypotheses, and it sharpens rather than weakens the contribution. Decoupling was expected to protect diversity (GA most stable, H1) and tight coupling was expected to be the liability under deception and multimodality (PSO most stressed, H2). The data show almost the opposite where it matters most: GA fails outright on deception and only partially handles non-stationarity, while the coupled architectures sustain diversity and solve every maze.

    The hypothesis tests confirm this reading. H1 is contradicted wherever GA actually searches over time, since on the trap and the wall it is the least diverse architecture rather than the most; the one maze where GA records high diversity is Parallel Paths, but there it halts in its first generation, so that value reflects the initial random spread and cannot support H1. H2 is contradicted on the trap and the wall, where GA rather than PSO or ACO is the worst case, and is supported only for PSO on multimodality, where social pressure does collapse the swarm onto a single route. H3 is also largely contradicted, since scale raises the diversity floor for every architecture and helps PSO rather than amplifying its weakness, although it does help GA most on the one maze where extra exploration can pay off.

    GA's complete failure on the trap is surprising enough to warrant ruling out an implementation bug, and the experiment itself does so. The same implementation solves Parallel Paths completely and the Sudden Wall partially, returning valid paths, so the encoding, operators, and goal test are correct and the failure is specific to deception rather than a general defect. On the trap, moreover, GA exhausts the full iteration budget within which ACO and PSO succeed and still reaches 0\% at every size, so the bottleneck is neither compute nor exploration.

    The mechanism behind this is shared memory. GA's decoupling is not a source of protective diversity but of wasted diversity: its agents explore independently, but with no shared memory the population cannot retain and build on the discovery that the dead end must be left. Selection on a distance-to-goal fitness actively pulls the population into the trap, and recombination of trapped parents produces trapped children, so the high apparent diversity is noise around a bad optimum rather than coverage of useful alternatives. The coupled architectures succeed for the same reason in reverse, because PSO's social term and ACO's pheromone are shared memory that lets a few lucky discoveries of the escape route propagate to the whole population. ACO is especially robust because its memory lives in the environment and decays, so an invalidated path is forgotten automatically, which is exactly why its Sudden-Wall recovery is near-instant rather than slow as predicted. GA's partial success on the wall fits the same account: lacking shared memory it cannot propagate the open route, but a large enough population occasionally contains a lineage that wanders onto it, which is why its wall success rises with scale while its trap success stays at zero. The right axis, then, is not that coupling hurts diversity but that shared memory converts exploration into exploitable progress, and its absence leaves a population unable to commit even when it has already found the answer.

    \subsection{Limitations}
    The results are conditional on fixed, untuned parameters. Per-maze tuning was avoided deliberately because it would over-fit and obscure the architectural differences under study, but it means the quantitative results depend on these parameter choices. A heavily tuned GA might do better than the untuned one shown here, though it exhausts the full iteration budget on the trap without a single success, which suggests tuning alone would not rescue it there.

    The dependence is sharpest for ACO's evaporation rate $\rho=0.1$, which produces long pheromone memory. As flagged in the proposal, this choice amplifies the adaptation effect: a low $\rho$ makes pheromone persist, so ACO's near-instant adaptation is a conservative result, and a higher $\rho$ would only make recovery faster. The central claim about ACO's robustness therefore does not depend on a favourable $\rho$. This is worth stating explicitly because the opposite worry, that $\rho$ was chosen to flatter ACO, would be natural, and the evidence runs the other way.

    Two further limitations bound the generality of the claims. The parameters adapted from Shrestha et al.\ were tuned for a different navigation dynamics, not a cell-by-cell grid, so they are a reasonable but not optimal starting point. And no prior work in the reference set evaluates dynamic environments, so the Sudden-Wall results extend the comparative literature into territory without a validated baseline, and should be read as a first characterisation rather than a tuned benchmark.

    \subsection{Future work}
    Three extensions follow directly. An evaporation-rate sweep for ACO and an inertia/social-coefficient sweep for PSO would map how far the observed robustness depends on memory persistence, turning the qualitative claim into a quantitative boundary. A GA variant with shared memory, such as an island model with migration, would test the mechanism directly: if the failure is really the absence of shared memory, adding it should rescue GA on the trap. Finally, scaling to larger and procedurally varied mazes, and to more than one disruption per run, would test whether the near-instant ACO and PSO recovery survives repeated non-stationarity or is an artefact of a single, mild disruption.

    \section{Conclusion}
    The information architecture of a population-based pathfinder strongly shapes its adaptive behaviour, but largely along the opposite of the hypothesised axis. Decoupling, expected to preserve useful diversity, instead leaves GA unable to escape a deceptive dead end at all, even when it exhausts its full search budget, and able to re-route around a new wall only partially and only with a large population. Shared memory, in the form of PSO's global best and especially ACO's decaying pheromone, is what lets a population convert exploration into a committed, recoverable solution, with the environmentally coupled architecture adapting the fastest to change. Increasing population size helps every architecture by raising the diversity floor, but it does not rescue GA on the deceptive trap. The line of work is therefore worth pursuing: a single controlled coupling axis already produces a clear, mechanistic, and partly counter-intuitive story, and the obvious next step, adding shared memory to GA and sweeping the memory parameters of PSO and ACO, is well-posed and likely to be informative.

    \newpage

    \section*{Author Contributions}
    \noindent\textbf{Manon Leijser:} Implemented the base mazes and the base GA algorithm; contributed to the project proposal and the renewed proposal; incorporated feedback on the renewed proposal; prepared and delivered the flash talk.

    \noindent\textbf{Nemo Ingendaa:} \textit{[contribution to be added]}

    \noindent\textbf{Robert Blaauwendraad:} Implemented the PSO and ACO algorithms and the path-encoded GA, redesigned the three maze environments, and built the experiment runner, analysis, and figure scripts; contributed to the writing of the project proposal and final report.

    \section*{Supplements}
    Source code, all maze-generation and experiment scripts, the full \texttt{runs.csv} (2{,}700 runs), figures, and a README with installation instructions, package/version pinning (\texttt{uv.lock}), a usage manual, and a sample run are available at:\\
    \textbf{GitHub:} \texttt{https://github.com/manon-ru/natcomproject}

    \bibliographystyle{plain}
    \bibliography{bibliography-project}

\end{document}